\documentclass[12 pt, letterpaper]{hmcpset}
\usepackage{hw-preamble}

\name{Jayden Thadani}
\class{Linear algebra done right}
\assignment{Exercises 7A --- Self-Adjoint and Normal Operators}
\duedate{}

\begin{document}

\begin{problem}[1]
	Suppose $n$ is a positive integer. Define $T \in \mathcal{L}(\mathbb{F}^{n})$ by
	\[
		T(z_{1}, \dots, z_{n}) = (0, z_{1}, \dots, z_{n - 1}).
	\]
	Find a formula for $T^{*}(z_{1}, \dots, z_{n})$.
\end{problem}

\begin{solution}
	\[
		\boxed{
			T^{*}(z_{1}, \dots, z_{n}) = (z_{2}, \dots, z_{n}, 0).
		}
	\]

	By definition, $\left < T(w_{1}, \dots, w_{n}), (z_{1}, \dots, z_{n}) \right > = \left < (w_{1}, \dots, w_{n}), T^{*}(z_{1}, \dots, z_{n}) \right >$. However, we know already that
	\[
		\begin{split}
			\left < T(w_{1}, \dots, w_{n}), (z_{1}, \dots, z_{n}) \right > & = \left < (0, w_{1}, \dots, w_{n - 1}), (z_{1}, \dots, z_{n}) \right > \\
										       & = w_{1} \overline{z_{2}}, + \dots + w_{n - 1} \overline{z_{n}} \\
										       & = \left < (w_{1}, \dots, w_{n}), (z_{2}, \dots, z_{n}, 0) \right >.
		\end{split}
	\]
	Thus, we must have $T^{*}(z_{1}, \dots, z_{n}) = (z_{2}, \dots, z_{n}, 0)$.
\end{solution}

\pagebreak

\begin{problem}[2]
	Suppose $T \in \mathcal{L}(V, W)$. Prove that
	\[
		T = 0 \iff T^{*} = 0 \iff T^{*}T = 0 \iff T T^{*} = 0.
	\]
\end{problem}

\begin{solution}
	Suppose $T = 0$. Then $\left < Tv, w \right > = 0$ for all $v, w \in V$, and thus, $\left < T^{*} w, v \right > = 0$ for all $v, w \in V$, including $v = T^{*} w$. That is, $\lVert T^{*} w \rVert = 0$ for all $w \in V$, and thus, $T^{*} = 0$. Since $(T^{*})^{*} = T$, this also gives us the converse --- that $T^{*} = 0 \implies T = 0$.

	It's obvious that either $T = 0$ or $T^{*} = 0$ is sufficient to show that $T T^{*} = T^{*} T = 0$. Now we will show the opposite direction. Suppose $T T^{*} = 0$. Then $\left < T T^{*} v, w \right > = 0$ always. However $\left < T T^{*} v, w \right > = \left < T^{*} v, T^{*} w \right >$ and by choosing $v = w$, we get that $\lVert T^{*} v \rVert = 0$ always. That is, $T^{*} = 0$, which gives us $T^{*} T = 0$. Thus, we have $T T^{*} = 0 \implies T^{*} = 0 \implies T^{*} T = 0$ (and also implies $T = 0$). Similarly, if  $T^{*} T = 0$, we have $\left < T^{*} T v, w \right > = \left < Tv, Tw \right >$ as $0$ always. Thus, $T^{*} T = 0 \implies T = 0 \implies T T^{*} = 0$ (and also implies $T^{*} = 0$). This completes the equivalence.
\end{solution}

\pagebreak

\begin{problem}[3]
	Suppose $T \in \mathcal{L}(V)$ and $\lambda \in \mathbb{F}$. Prove that
	\[
		\lambda \text{ is an eigenvalue of } T \iff \overline{\lambda} \text{ is an eigenvalue of } T^{*}.
	\]
\end{problem}

\begin{solution}
	Suppose $\lambda$ is an eigenvalue of $T$ with eigenvector $v$. This is implies that $T - \lambda I$ is non-invertible, and thus, so is its adjoint, $T^{*} - \overline{\lambda} I$. Since $T^{*} - \overline{\lambda} I$ is non-invertible, $T^{*}$ has $\overline{\lambda}$ as an eigenvalue. Since $(T^{*})^{*} = T$, the converse is also implied.
\end{solution}

\pagebreak

\begin{problem}[4]
	Suppose $T \in \mathcal{L}(V)$ and $U$ is a subspace of $V$. Prove that
	\[
		U \text{ is invariant under } T \iff U^{\perp} \text{ is invariant under } T^{*}.
	\]
\end{problem}

\begin{solution}
	First suppose $U$ is invariant under $T$. Then for any $u \in U$ and $w \in U^{\perp}$, we have $Tu \in U$ as well, and thus, $\left < Tu, w  \right > = 0$. Thus, $\left < u, T^{*}w \right > = 0$ for any $w \in U^{\perp}$, $u \in U$. For this, we must have $T^{*} w \in U^{\perp}$ for any $w \in U^{\perp}$. That is, $U^{\perp}$ is invariant under $T$. If $V$ is finite dimensional, the converse follows just from the $-^{*}$ and $-^{\perp}$ operators being involutions.
\end{solution}

\pagebreak

\begin{problem}[5]
	Suppose $T \in \mathcal{L}(V, W)$. Suppose $e_{1}, \dots, e_{n}$ is an orthonormal basis of $V$ and $f_{1}, \dots, f_{m}$ is an orthonormal basis of $W$. Prove that
	\[
		\lVert T e_{1} \rVert^{2} + \dots + \lVert T e_{n} \rVert^{2} = \lVert T^{*} f_{1} \rVert^{2} + \dots + \lVert T^{*} f_{m} \rVert^{2}.
	\]
\end{problem}

\begin{solution}
	Since $e_{1}, \dots, e_{n}$ and $f_{1}, \dots, f_{m}$ are orthonormal bases we can write each $T e_{j}$ in the basis $f_{1}, \dots, f_{m}$ and each $T^{*} f_{k}$ in the basis $e_{1}, \dots, e_{n}$. Thus, we can write
	\[
		\begin{split}
			\lVert T e_{1} \rVert^{2} + \dots + \lVert T e_{n} \rVert^{2} & = \sum_{j = 1}^{n} \left \Vert \sum_{k = 1}^{m} \left < T e_{j}, f_{k} \right > f_{k} \right \Vert^{2} \\
										      & = \sum_{k = 1}^{m} \sum_{j = 1}^{n} \lvert \left < T e_{j}, f_{k} \right > \rvert^{2} \\
										      & = \sum_{j = 1}^{n} \sum_{k = 1}^{m} \lvert \left < e_{j}, T^{*} f_{k} \right > \rvert^{2} \\
										      & = \sum_{k = 1}^{m} \sum_{j = 1}^{n} \lvert \left < T^{*} f_{k}, e_{j} \right > \rvert^{2} \\
										      & = \sum_{k = 1}^{m} \left \lVert \sum_{j = 1}^{n} \left < T^{*} f_{k}, e_{j} \right > e_{j} \right \rVert^{2} \\
										      & = \lVert T^{*} f_{1} \rVert^{2} + \dots + \lVert T^{*} f_{m} \rVert^{2}.
		\end{split}
	\]
\end{solution}

\pagebreak

\begin{problem}[6]
	Suppose $T \in \mathcal{L}(V, W)$. Prove that
	\begin{enumerate}
		\item $T$ is injective $\iff$ $T^{*}$ is surjective;
		\item $T$ is surjective $\iff$ $T^{*}$ is injective.
	\end{enumerate}
\end{problem}

\begin{solution}
	Since $\dim \operatorname{range} T^{*} = \dim (\operatorname{null} T)^{\perp}$ and $V, W$ are finite-dimensional, we have $\dim \operatorname{range} T^{*} = \dim V - \dim \operatorname{null} T$. Thus, $T$ is injective if and only if $T^{*}$ is surjective. The second equivalence follows from $-^{*}$ being an involution.
\end{solution}

\pagebreak

\begin{problem}[7]
	Prove that if $T \in \mathcal{L}(V, W)$, then
	\begin{enumerate}
		\item $\dim \operatorname{null} T^{*} = \dim \operatorname{null} T + \dim W - \dim V$;
		\item $\dim \operatorname{range} T^{*} = \dim \operatorname{range} T$.
	\end{enumerate}
\end{problem}

\begin{solution}
	Note again that $\dim \operatorname{range} T^{*} = \dim (\operatorname{null} T)^{\perp} = \dim V - \dim \operatorname{null} T$ and by symmetry, $\dim \operatorname{range} T = \dim W - \dim \operatorname{null} T^{*}$. Also, the rank-nullity theorem gives us that $\dim \operatorname{null} T + \dim \operatorname{range} T = \dim V$, and similarly $\dim \operatorname{null} T^{*} + \dim \operatorname{range} T^{*} = \dim W$
	\begin{enumerate}
		\item From the equalities above, it is only a matter of algebraic manipulation to show that
			\[
				\begin{split}
					\dim \operatorname{null} T^{*} & = \dim V - \dim \operatorname{range} T^{*} \\
								       & = \dim W - (\dim V - \dim \operatorname{null} T) \\
								       & = \dim \operatorname{null} T + \dim W - \dim \operatorname{null} T.
				\end{split}
			\]
		\item Similarly, it is easily shown that
			\[
				\begin{split}
					\dim \operatorname{range} T^{*} & = \dim V - \dim \operatorname{null} T \\
									& = \dim \operatorname{range} T.
				\end{split}
			\]
	\end{enumerate}
\end{solution}

\pagebreak

\begin{problem}[8]
	Suppose $A$ is an $m$-by-$n$ matrix with entries in $\mathbb{F}$. Use (b) in Exercise 7 to prove that the row rank of $A$ equals the column rank of $A$.
\end{problem}

\begin{solution}
	The row rank of an $m$-by-$n$ matrix in $\mathbb{F}$ is just the same as the column rank of its transpose. We claim that this is also the same as the column rank of its conjugate transpose.

	Suppose we have a matrix $A$ with transpose $A^{\top}$ and conjugate transpose $A^{*}$. Notice that $a_{i_{1}} A^{\top}_{j, 1} + \dots + a_{i_{k}} A^{\top}_{j, i_{k}} = 0$ if and only if $\overline{a_{i_{1}}} A^{*}_{j, i_{1}} + \dots + \overline{a_{i_{k}}} A^{*}_{j, i_{k}} = 0$. Thus, any $k$ columns of  $A^{\top}$ are linearly independent if and only if the same $k$ columns of  $A^{*}$ are linearly independent. That is $A^{\top}$ and $A^{*}$ have the same column rank, and thus, the column rank of $A^{*}$ is the row rank of $A$.

	Note $\operatorname{range} T^{*}$ is the column rank of $A^{*}$, and thus, the row rank of $A$. Since, the column rank of $A$ is $\operatorname{range} T = \operatorname{range} T^{*}$, the column rank of $A$ is the row rank of $A$.
\end{solution}

\pagebreak

\begin{problem}[9]
	Prove that the product of two self-adjoint operators on $V$ is self-adjoint if and only if the two operators commute.
\end{problem}

\pagebreak

\begin{problem}[10]
	Suppose $\mathbb{F} = \mathbb{C}$ and $\mathrm{T} \in \mathcal{L}(V)$. Prove that $T$ is self-adjoint if and only if
	 \[
		\left < Tv, v \right > = \left < T^{*}v, v \right >
	\]
	for all $v \in V$.
\end{problem}

\pagebreak

\begin{problem}[11]
	Define an operator $S: \mathbb{F}^{2} \to \mathbb{F}^{2}$ by $S(w, z) = (-z, w)$.
	\begin{enumerate}
		\item Find a formula for $S^{*}$.
		\item Show that $S$ is normal but not self-adjoint.
		\item Find all eigenvalues of $S$.
	\end{enumerate}
\end{problem}

\pagebreak

\begin{problem}[12]
	An operator $B \in \mathcal{L}(V)$ is called \emph{skew} if
	\[
		B^{*} = -B.
	\]
	Suppose that $T \in \mathcal{L}(V)$. Prove that $T$ is normal if and only if there exist commuting operators $A$ and $B$ such that $A$ is self-adjoint, $B$ is a skew operator, and $T = A + B$.
\end{problem}

\pagebreak

\begin{problem}[13]
	Suppose $\mathbb{F} = \mathbb{R}$. Define $\mathcal{A} \in \mathcal{L}(\mathcal{L}(V))$ by $\mathcal{A} T = T^{*}$ for all $T \in \mathcal{L}(V)$.
	\begin{enumerate}
		\item Find all eigenvalues of $\mathcal{A}$.
		\item Find the minimal polynomial of $\mathcal{A}$.
	\end{enumerate}
\end{problem}

\pagebreak

\begin{problem}[14]
	Define an inner product on $\mathcal{P}_{2}(\mathbb{R})$ by $\left < p, q \right > = \int_{0}^{1} pq$. Define an operator $T \in \mathcal{L}(\mathcal{P}_{2}(\mathbb{R}))$ by
	\[
		T(ax^{2} + bx + c) = bx
	\]
	\begin{enumerate}
		\item Show that this inner product, the operator is not self-adjoint.
		\item The matrix of $T$ with respect to the basis $1, x, x^{2}$ is
			\[
				\begin{bmatrix}
					0 & 0 & 0 \\
					0 & 1 & 0 \\
					0 & 0 & 0
				\end{bmatrix}
			\]
			This matrix equals its conjugate transpose, even though $T$ is not self-adjoint. Explain why this is not a contradiction.
	\end{enumerate}
\end{problem}

\pagebreak

\begin{problem}[15]
	Suppose $T \in \mathcal{L}(V)$ is invertible. Prove that
	\begin{enumerate}
		\item $T$ is self-adjoint $\iff$ $T^{-1}$ is self-adjoint;
		\item $T$ is normal $\iff$ $T^{-1}$ is normal.
	\end{enumerate}
\end{problem}

\pagebreak

\begin{problem}[16]
	Suppose $\mathbb{F} = \mathbb{R}$.
	\begin{enumerate}
		\item Show that the set of self-adjoint operators on $V$ is a subspace of $\mathcal{L}(V)$.
		\item What is the dimension of the subspace of $\mathcal{L}(V)$ in (a) (in terms of $\dim V$)?
	\end{enumerate}
\end{problem}

\pagebreak

\begin{problem}[17]
	Suppose $\mathbb{F} = \mathbb{C}$. Show that the set of self-adjoint operators on $V$ is not a subspace of $\mathcal{L}(V)$.
\end{problem}

\pagebreak

\begin{problem}[18]
	Suppose $\dim V \le 2$. Show that the set of normal operators on $V$ is not a subspace of $\mathcal{L}(V)$.
\end{problem}

\pagebreak

\begin{problem}[19]
	Suppose $T \in \mathcal{L}(V)$ and $\lVert T^{*} v \rVert \le \lVert Tv \rVert$ for every $v \in V$. Prove that $T$ is normal.
\end{problem}

\pagebreak

\begin{problem}[20]
	Suppose $P \in \mathcal{L}(V)$ is such that $P^{2} = P$. Prove that the following are equivalent.
	\begin{enumerate}
		\item $P$ is self adjoint.
		\item $P$ is normal.
		\item There is a subspace of $U$ of $V$ such that $P = P_{U}$.
	\end{enumerate}
\end{problem}

\pagebreak

\begin{problem}[21]
	Suppose $D: \mathcal{P}_{8}(\mathbb{R}) \to \mathcal{P}_{8}(\mathbb{R})$ is the differentiation operator defined by $Dp = p'$. Prove that there does not exist an inner product on $\mathcal{P}_{8}(\mathbb{R})$ that makes $D$ a normal operator.
\end{problem}

\pagebreak

\begin{problem}[22]
	Give an example of an operator $T \in \mathcal{L}(\mathbb{R}^{3})$ such that $T$ is normal but not self-adjoint.
\end{problem}

\pagebreak

\begin{problem}[23]
	Suppose $T$ is a normal operator on $V$. Suppose also that $v, w \in V$ satisfy the equations
	\[
		\lVert v \rVert = \lVert w \rVert = 2, \quad Tv = 3, \quad Tw = 4w.
	\]
	Show that $\lVert T(v + w) \rVert = 10$.
\end{problem}

\pagebreak

\begin{problem}[24]
	Suppose $T \in \mathcal{L}(V)$ and
	\[
		a_{0} + a_{1} z + a_{2} z^{2} + \dots + a_{m - 1} z^{m - 1} + z^{m}
	\]
	is the minimal polynomial of $T$. Prove that the minimal polynomail of $T^{*}$ is
	\[
		\overline{a_{0}} + \overline{a_{1}} z + \overline{a_{2}} z^{2} + \dots + \overline{a_{m - 1}} z^{m - 1} + z^{m}.
	\]
\end{problem}

\pagebreak

\begin{problem}[25]
	Suppose $T \in \mathcal{L}(V)$. Prove that $T$ is diagonalisable if and only if $T^{*}$ is diagonalisable.
\end{problem}

\pagebreak

\begin{problem}[26]
	Fix $u, x \in V$. Define $T \in \mathcal{L}(V)$ by $Tv = \left < v, u \right > x$ for every $v \in V$.
	\begin{enumerate}
		\item Prove that if $V$ is a real vector space, then $T$ is self-adjoint if and only if the list  $u, x$ is linearly dependent.
		\item Prove that $T$ is normal if and only if the list $u, x$ is linearly dependent.
	\end{enumerate}
\end{problem}

\pagebreak

\begin{problem}[27]
	Suppose $T \in \mathcal{L}(V)$ is normal. Prove that
	\[
		\operatorname{null} T^{k} = \operatorname{null} T \text{ and } \operatorname{range} T^{k} = \operatorname{range} T
	\]
	for every positive integer $k$.
\end{problem}

\pagebreak

\begin{problem}[28]
	Suppose $T \in \mathcal{L}(V)$ is normal. Prove that if $\lambda \in \mathbb{F}$, then the minimal polynomial of $T$ is not a polynomial multiple of $(x - \lambda)^{2}$.
\end{problem}

\pagebreak

\begin{problem}[29]
	Prove or give a counterexample: If $T \in \mathcal{L}(V)$ and there is an orthonormal basis $e_{1}, \dots, e_{n}$ of $V$ such that $\lVert T e_{k} \rVert = \lVert T^{*} e_{k} \rVert$ for each $k = 1, \dots, n$, then $T$ is normal.
\end{problem}

\pagebreak

\begin{problem}[30]
	Suppose $T \in \mathcal{L}(\mathbb{F}^{3})$ is normal and $T(1, 1, 1) = (2, 2, 2)$. Suppose $(z_{1}, z_{2}, z_{3}) \in \operatorname{null} T$. Prove that $z_{1} + z_{2} + z_{3} = 0$.
\end{problem}

\pagebreak

\begin{problem}[31]
	Fix a positive integer $n$. In the inner product space of continuous real-valued functions on  $[-\pi, \pi]$ with inner product $\left < f, g \right > = \int_{- \pi}^{\pi} fg$, let
	\[
		V = \operatorname{span}(1, \cos{x}, \cos{2x}, \dots, \cos{nx}, \sin{x}, \sin{2x}, \dots, \sin{nx}).
	\]
	\begin{enumerate}
		\item Define $D \in \mathcal{L}(V)$ by $Df = f'$. Show that $D^{*} = - D$. Conclude that $D$ is normal but not self-adjoint.
		\item Define $T \in \mathcal{L}(V)$ by $Tf = f''$. Show that $T$ is self-adjoint.
	\end{enumerate}
\end{problem}

\pagebreak

\begin{problem}[32]
	Suppose $T : V \to W$ is a linear map. Show that under the standard identification of $V$ with $V^{\vee}$ and the corresponding identification of $W$ with $W^{\vee}$ the adjoint map $T^{*} : W \to V$ corresponds to the dual map $T^{\vee} : W^{\vee} \to V^{\vee}$. More precisely, show that
	\[
		T^{\vee}(\varphi_{w}) = \varphi_{T^{*} w}
	\]
	for all $w \in W$.
\end{problem}

\end{document}
